\section{Results}
\label{sec:results}

The results are organized around the three benchmark tiers defined in
Section~\ref{sec:benchmark_design}. The point of that structure is to separate three questions
that the earlier draft treated together: how $L_0$ behaves on a tractable comparison, how the
methods degrade as the target system grows, and which methods remain usable near the production
problem. Numerical entries are still marked as placeholders pending completed benchmark runs.

\subsection{Tier 1: tractable benchmark comparisons}

Table~\ref{tab:comparison} reports the tractable benchmark. The mixed-target setting compares $L_0$
with GREG on the same exported calibration package. The count-style setting compares $L_0$ with IPF
on an IPF-eligible subset derived from that same package.

\begin{table}[ht]
\centering
{\tablefont
\begin{tabular}{lrrrr}
\toprule
Method & Mean ARE (\%) & Median ARE (\%) & Runtime & Sparsity control \\
\midrule
$L_0$ (this paper) & \tbc & \tbc & \tbc & Yes (configurable $\lambda_{L_0}$) \\
IPF & \tbc & \tbc & \tbc & No \\
GREG & \tbc & \tbc & \tbc & No \\
\bottomrule
\end{tabular}
}
\caption{Comparison of $L_0$ regularization with iterative proportional fitting (IPF) and generalized regression (GREG) estimators on the same calibration matrix and target set.}
\label{tab:comparison}
\tablenote{Note: \tbc[Results to be populated after running IPF and GREG baselines. ARE = absolute relative error across all achievable targets.]}
\end{table}


The tractable benchmark is the most direct test of whether $L_0$ performs well only because it was
designed for the full production setting. In the mixed-target benchmark, the key comparisons are
accuracy, runtime, and whether GREG produces negative weights. In the count-style benchmark, the
key question is whether $L_0$ remains competitive with IPF even when the target system is limited
to the class of margins that IPF handles naturally.

\subsection{Tier 2: scaling frontier}

Table~\ref{tab:scaling_frontier} summarizes the scaling ladder. Each rung uses the same basic
cloned unit universe while increasing the number of benchmark targets. This isolates how method
performance changes with constraint volume rather than with a changing data-preparation pipeline.

\begin{table}[ht]
\centering
{\tablefont
\resizebox{\textwidth}{!}{%
\begin{tabular}{rrrrll}
\toprule
Tier & Targets & $L_0$ runtime & GREG runtime & IPF runtime & Main outcome \\
\midrule
1 & 250 & \tbc & \tbc & \tbc & \tbc \\
2 & 500 & \tbc & \tbc & \tbc & \tbc \\
3 & 1{,}000 & \tbc & \tbc & \tbc & \tbc \\
4 & 2{,}500 & \tbc & \tbc & \tbc & \tbc \\
5 & 5{,}000 & \tbc & \tbc & \tbc & \tbc \\
6 & 10{,}000 & \tbc & \tbc & \tbc & \tbc \\
7 & \tbc[full selected tier] & \tbc & \tbc & \tbc & \tbc \\
\bottomrule
\end{tabular}
}
}
\caption{Tier 2 scaling frontier. Each row uses the same cloned unit universe and increases the
number of selected targets. Runtime cells should be replaced by failure labels when a method does
not complete.}
\label{tab:scaling_frontier}
\tablenote{The final paper may split this table into separate runtime and failure-mode tables if
the notes column becomes too dense.}
\end{table}


The scaling frontier is expected to show more than runtime growth. It should also identify where a
method becomes numerically unstable, fails to converge, or stops being interpretable as a credible
comparison. For GREG, the relevant failure modes include memory pressure, unstable linear algebra,
and negative weights. For IPF, the relevant limits include inapplicability outside count-style
margins and non-convergence when the margin system becomes too hard to reconcile.

\subsection{Tier 3: production feasibility}

Table~\ref{tab:production_feasibility} reports the production-feasibility case. This table is
designed to be concrete. It should state which method completed, which method failed, how long the
run took, and what the failure looked like when a method did not finish.

\begin{table}[ht]
\centering
{\tablefont
\resizebox{\textwidth}{!}{%
\begin{tabular}{p{0.16\textwidth}p{0.17\textwidth}p{0.17\textwidth}rrp{0.19\textwidth}}
\toprule
Method & Target set & Status & Runtime & Peak memory & Notes \\
\midrule
$L_0$ & Full production mixed targets & \tbc & \tbc & \tbc & Retained records: \tbc \\
GREG & Full production mixed targets & \tbc & \tbc & \tbc & Negative weights / solver status: \tbc \\
IPF & Production count-style subset & \tbc & \tbc & \tbc & Iterations or failure mode: \tbc \\
\bottomrule
\end{tabular}
}
}
\caption{Tier 3 production-feasibility benchmark. This table should report completion and failure
states as directly as possible.}
\label{tab:production_feasibility}
\tablenote{The IPF row is reported on a count-style production subset because IPF is not used as a
general solver for the full mixed-target system.}
\end{table}


This tier matters because a tractable benchmark alone does not establish practical usefulness. A
method can look reasonable at a few hundred targets and still fail to support a deployable
subnational dataset. The production-feasibility case links the benchmark back to the actual
\texttt{policyengine-us-data} use case.

\subsection{Supporting diagnostics}

The tiered benchmark tables provide the main narrative. The following diagnostics support that
narrative by showing how the fitted $L_0$ solution behaves within a completed run.

\subsubsection{Accuracy by geography}

Table~\ref{tab:calibration_accuracy} reports absolute relative error by geographic level for the
selected benchmark or production configuration.

\begin{table}[ht]
\centering
{\tablefont
\begin{tabular}{lrrrr}
\toprule
Geographic level & Targets & Mean ARE (\%) & Median ARE (\%) & Max ARE (\%) \\
\midrule
National & \tbc & \tbc & \tbc & \tbc \\
State & \tbc & \tbc & \tbc & \tbc \\
Congressional district & \tbc & \tbc & \tbc & \tbc \\
\midrule
\textbf{All achievable} & \tbc & \tbc & \tbc & \tbc \\
\bottomrule
\end{tabular}
}
\caption{Calibration accuracy by geographic level, measured as absolute relative error (ARE) across achievable targets under the local preset ($\lambda_{L_0} = 10^{-8}$).}
\label{tab:calibration_accuracy}
\tablenote{Note: \tbc[Values to be populated from \texttt{unified\_diagnostics.csv} after a completed pipeline run.]}
\end{table}


The benchmark manifests determine the exact number of targets in each row. Those counts should come
from the exported package rather than from the full target database used by the pipeline, which can
include multiple years or inactive target families that are not part of a given experiment.

\subsubsection{Sparsity and convergence}

The sparsity penalty produces exact zeros in the retained weight vector, so dataset size is an
explicit output of the optimization rather than a post-processing step. For the settings reported
here, the main summary objects are the retained-record count, the effective sample size, and the
distribution of fitted weights.

Figure~\ref{fig:convergence} plots mean absolute relative error over training epochs for the
selected $L_0$ runs.

\begin{figure}[ht]
    \centering
    \fbox{\parbox{0.9\textwidth}{\centering\textit{[Convergence curves to be generated from
    calibration logs]}}}
    \caption{Convergence diagnostics for representative $L_0$ runs. The final figure will report
    the calibration loss over epochs for the tractable benchmark and the production-feasibility
    run.}
    \label{fig:convergence}
\end{figure}

These diagnostics are secondary to the benchmark tiers, but they are still important. They show
whether a feasible production run is achieving that feasibility by keeping errors low, by keeping
weights stable, and by pruning the cloned universe in a controlled way rather than through ad hoc
thresholding.
